\chapter{Минимальная динамическая parent-архитектура charged mediator}
\label{ch:v8-charged-mediator-minimal-dynamic-parent}

\section{Общая ячейка требует ровно две новые линии}

Существующая тёплиц-ячейка полного шума имеет вид
\begin{equation}
 \mathcal K_{43}=\mathbb C|0\rangle\oplus\mathcal N_{42}.
 \label{eq:v8-mediator-dynamic-old-cell}
\end{equation}
Для нового бинарного перехода одна комплексная charged-линия достаточна,
но не замкнута относительно Real-структуры. Поэтому минимальное расширение
равно
\begin{equation}
 \mathcal K_{45}=\mathbb C|0\rangle\oplus\mathcal N_{42}
 \oplus\mathbb C|+\rangle\oplus\mathbb C|-\rangle,
 \qquad \dim\mathcal K_{45}=45.
 \label{eq:v8-mediator-dynamic-extended-cell}
\end{equation}
На двух новых линиях
\begin{equation}
 Y_E|\pm\rangle=\pm|\pm\rangle,
 \qquad J_E|+\rangle=|-\rangle,
 \qquad H_E|\pm\rangle=7\gamma|\pm\rangle.
 \label{eq:v8-mediator-dynamic-real-energy-pair}
\end{equation}
Удаление любой из них разрушает Real-замыкание. Следовательно, переход
$43\to45$ минимален среди расширений общей ячейки.

\section{Локальный энергосохраняющий collision-parent}

На системном Real-дубле используем базис
$|u\rangle,|d\rangle,|u^c\rangle,|d^c\rangle$ и данные
\begin{equation}
 Y_S=\operatorname{diag}(1,0,-1,0),
 \qquad H_S=7\gamma\operatorname{diag}(1,0,1,0).
 \label{eq:v8-mediator-dynamic-system-data}
\end{equation}
Положим
\begin{equation}
 L_-=|d\rangle\langle u|,
 \qquad L_+=|d^c\rangle\langle u^c|,
 \qquad b_\pm=|\pm\rangle\langle0|.
 \label{eq:v8-mediator-dynamic-jumps}
\end{equation}
Минимальное Real-замкнутое взаимодействие есть
\begin{equation}
 G_g=g\left(L_-\otimes b_++L_+\otimes b_-+\mathrm{h.c.}\right).
 \label{eq:v8-mediator-dynamic-interaction}
\end{equation}
Оно удовлетворяет двум точным законам сохранения:
\begin{equation}
 [Y_S\otimes1+1\otimes Y_E,G_g]=0,
 \qquad [H_S\otimes1+1\otimes H_E,G_g]=0.
 \label{eq:v8-mediator-dynamic-conservation}
\end{equation}
После удаления множителя $g$ оператор имеет ранг четыре и спектр
\begin{equation}
 \Spec(G_g/g)=\{-1,-1,0^{\times8},1,1\}.
 \label{eq:v8-mediator-dynamic-spectrum}
\end{equation}

\section{Ортогональность нового шумового блока}

Матричные элементы из нейтрального вакуума точно возвращают оба
сопряжённых перехода:
\begin{equation}
 \langle+|G_g|0\rangle=gL_-,
 \qquad \langle-|G_g|0\rangle=gL_+.
 \label{eq:v8-mediator-dynamic-extraction}
\end{equation}
Их cross-covariance исчезает, а первый и второй моменты равны
\begin{equation}
 \langle0|G_g|0\rangle=0,
 \qquad
 \langle0|G_g^2|0\rangle
 =g^2(L_-^\dagger L_-+L_+^\dagger L_+).
 \label{eq:v8-mediator-dynamic-moments}
\end{equation}
Поэтому новый блок присоединяется к старому 42-frame ортогональной прямой
суммой. В слабом collision-пределе получаем
\begin{equation}
 \mathcal L_{44}=\mathcal L_{42}
 +g^2\mathcal D_{L_-}+g^2\mathcal D_{L_+}.
 \label{eq:v8-mediator-dynamic-generator}
\end{equation}
Если $\mathcal L_{42}$ уже примитивен, то добавление неотрицательных
дирихле-форм не расширяет его неподвижную алгебру:
\begin{equation}
 \operatorname{Fix}(\mathcal L_{44})
 \subseteq\operatorname{Fix}(\mathcal L_{42})=\mathbb C1.
 \label{eq:v8-mediator-dynamic-fixed-algebra}
\end{equation}

\section{Вакуум имеет локального parent, транспорт остаётся Флоке}

На каждой 45-мерной ячейке положительный оператор
\begin{equation}
 h_{\rm cell}=\sum_{a=1}^{42}|a\rangle\langle a|
 +7\bigl(|+\rangle\langle+|+|-\rangle\langle-|\bigr)
 \label{eq:v8-mediator-dynamic-cell-parent}
\end{equation}
имеет единственное нулевое состояние $|0\rangle$. Сумма его трансляций
задаёт frustration-free parent тензорного вакуума. Один общий такт
\begin{equation}
 V_{45}=(I\otimes S_{45})U_{\rm col}^{(0)}
 \label{eq:v8-mediator-dynamic-floquet-step}
\end{equation}
точно повторяет расширенный канал на свежих ячейках. Однако
$\operatorname{ind}_{\rm GNVW}(S_{45})=45$, поэтому этот сдвиг не является
конечновременной эволюцией локального гамильтониана из компоненты
тождественного автоморфизма.

Итак, условная архитектура закрыта полностью:
\begin{equation}
 N_{\rm conditional\ architecture}=10/10.
 \label{eq:v8-mediator-dynamic-conditional-ledger}
\end{equation}
Но она использует пять невыведенных входов:
\begin{equation}
 N_{\rm inherited\ parent}=0/5:
 \{H_{23/24},\,7\gamma,\,g,\,S_{45},\,\tau_C\}.
 \label{eq:v8-mediator-dynamic-parent-ledger}
\end{equation}

\begin{theorem}[Минимальная 45-мерная динамическая архитектура]
Существующий 43-мерный тёплиц-носитель допускает минимальное Real-замкнутое
расширение двумя charged-линиями. Оно даёт точный gauge- и
energy-сохраняющий collision-parent, ортогонально расширяет генератор с 42
до 44 каналов и имеет локальный frustration-free parent тензорного
вакуума. Физические коэффициенты и транспортный закон при этом не
наследуются.
\end{theorem}

\begin{nogocriterion}[Архитектурная полнота не есть происхождение]
Нельзя считать число $45$ селектором щели $7\gamma$ или coupling $g$.
Локальный parent вакуума не порождает сдвиг $S_{45}$, а точный Флоке-шаг
не задаёт физическую длительность столкновения. Условные расширения
$H_{23/24}$ также нельзя объявлять уже выведенными из raw $H_{21}$.
\end{nogocriterion}

\begin{protocol}\begin{enumerate}
\item старая ячейка: \textbf{$1+42=43$};
\item минимально добавленных Real-линий: \textbf{$2$};
\item расширенная ячейка: \textbf{$1+42+2=45$};
\item полный заряд сохраняется: \textbf{да};
\item свободная энергия сохраняется: \textbf{да};
\item ранг нового interaction-блока: \textbf{$4$};
\item cross-covariance двух новых каналов: \textbf{$0$};
\item полный weak-limit frame: \textbf{$44$ канала};
\item локальный parent product-vacuum: \textbf{есть};
\item GNVW-индекс сдвига: \textbf{$45$};
\item условная динамическая архитектура: \textbf{$10/10$};
\item inherited physical parent: \textbf{$0/5$};
\item следующий гейт: \textbf{admission динамических parent-данных}.
\end{enumerate}\end{protocol}

Машинный сертификат сохранён в
\path{s2t_v8_baryon_c0_singlet_triplet_central_gap_isotypic_channel_extra_edge_mass_two_source_cubic_incidence_up_down_binary_selector_minimal_typed_discriminator_binary_carrier_charged_environment_mediator_minimal_dynamic_parent_architecture_gate_results.json}.

\begin{thebibliography}{9}
\bibitem{v8-mediator-dynamic-attal} S.~Attal and Y.~Pautrat,
\emph{From Repeated to Continuous Quantum Interactions},
Ann. Henri Poincar\'e 7 (2006) 59; arXiv:math-ph/0311002.
\bibitem{v8-mediator-dynamic-gnvw} D.~Gross, V.~Nesme, H.~Vogts and
R.~F.~Werner, \emph{Index Theory of One Dimensional Quantum Walks and
Cellular Automata}, Commun. Math. Phys. 310 (2012) 419;
arXiv:0910.3675.
\end{thebibliography}